% v21: five necessary support types, with sharper book/G5_15/local-D reductions.
% The sf20 aliases are retained for references in earlier modules.
\subsection{残る五型と各支持内の追加の必要条件}
\label{sf21:section}\label{sf20:section}

\begin{corollary}[非空支持の連結性と二価頂点の制約]
\label{sf21:no-leaves}\label{sf20:no-leaves}
正の局所整数 charts，marked puncture の整数 clutching，対向 $D$ の相殺，
および node 非負性を満たす，五次の標準 $K_5$ branch graph を考える．
高分母 support $H=\{ij:q_{ij}>1\}$ が非空なら，
$H$ は五頂点全てを使う二点連結 graph であり，最小次数は二以上である．
さらに
\[
 |E(H)|\ge7,\qquad b_1(H)\ge3
\]
であり，各二価頂点の二つの高近傍は，ともに次数四を持つ．
特に二価頂点は高三角形に属する．
高辺が七本なら $H$ は book 型 $G_{5,13}$ である．
二価頂点の個数は零，一，三のいずれかであり，
一つの場合は $G_{5,15}$，三つの場合は book 型である．
一つの core の高分母 ray が高々一本なら，全ての分母は一である．
\end{corollary}

\begin{proof}
系~\ref{kg17:low-core-propagation} により，非空の場合に孤立頂点はない．
木成分，単純閉路成分，二価二頂点を持つ三角形，葉付き四角形，
$K_{2,3}$，diamond，二葉三角形および四使用頂点 $K_4$ の障害を適用すると，
五頂点上の十一型が残る．
定理~\ref{hs18:no-house}，\ref{gten18:no-support}，\ref{lh18:no-K4-leaf} により，
$G_{5,5},G_{5,10},G_{5,12}$ が除かれる．
この段階で葉を持つ型，または六辺の型は $G_{5,7},G_{5,9}$ だけであり，
定理~\ref{g719:no-G5-7} と定理~\ref{dl19:no-G5-9} がそれらを排除する．
従って最小次数二以上，辺数七以上，$b_1(H)=|E(H)|-4\ge3$ を得る．
定理~\ref{g51120:no-support} により $G_{5,11}$ も除かれる．
次の五代表の表から，二価頂点の両近傍が次数四を持つことと，
七辺の場合および二価頂点数についての主張が従う．

二点連結性は最小次数と辺数からも直接従う．切断頂点 $v$ があるなら，
$H-v$ の各成分は最小次数二の条件により二頂点以上を持つ．
残る四頂点は二頂点ずつの二成分に分かれ，各成分の両頂点は互いに接続し，
ともに $v$ に接続しなければならない．
これは二つの三角形が $v$ だけを共有する六辺の graph であり，辺数七以上に反する．
最後の結論は，孤立頂点と葉がいずれも許されないことの言い換えである．
\end{proof}

\begin{corollary}[五型六十六標識への帰着]
\label{sf21:remaining-support}\label{sf20:remaining-support}
前系の仮定の下で，非空の高分母 support の必要候補は次の五型に限られる．
標識を保つ候補は六十六個である．$\overline H=K_5\setminus H$ とし，
補 graph の表示では孤立頂点を省略する．
\begin{center}\small
\begin{tabular}{crcrcr}
型&$|E|$&降順の次数列&$b_1$&$\overline H$&標識付き個数\\\hline
$G_{5,13}$&7&$(4,4,2,2,2)$&3&$K_3$&10\\
$G_{5,14}$&8&$(4,3,3,3,3)$&4&$2K_2$&15\\
$G_{5,15}$&8&$(4,4,3,3,2)$&4&$P_3$&30\\
$G_{5,16}$&9&$(4,4,4,3,3)$&5&$K_2$&10\\
$G_{5,17}$&10&$(4,4,4,4,4)$&6&$\varnothing$&1
\end{tabular}
\end{center}
識別子は固定し，欠番を詰めない．
これらは必要条件を通過した graph の個数であり，
整数接着データ，作用 atlas，または Lagrangian fibrations の存在数を表さない．
\end{corollary}

\begin{proof}
全 $2^{10}=1024$ 個の標識付き部分 graph に既述の障害を適用する．
$G_{5,11}$ は自己同型群の位数が二であるため，その排除は六十標識を除く．
五代表の自己同型群の位数は表の順に $12,8,4,12,120$ であり，
軌道数 $5!/|\operatorname{Aut}(H)|$ の和は
\[
 10+15+30+10+1=66
\]
となる．各代表の補 graph を取ると第五列が得られる．
\nolinkurl{verify_support_agenda_v20.py} は全 graph に対して，
固定した識別子と軌道数，各頂点の削除後の連結性，
二価頂点の二近傍の次数を確認する．二価頂点の確認は全標識を通して六十件となる．
結果は \nolinkurl{support_agenda_verification_v20.json} に全標識とともに保存する．
この有限 graph の列挙は，先行する整数接着の排除定理を適用するものであり，
それらの証明の代わりではない．
\end{proof}

\paragraph{Book に残る二つの sector．}
残る低辺は高々三本であり，三本なら低三角形を作る book 型だけである．
固定した一 core に対する有限完成法（定理~\ref{bk19:finite-completion}）と，
全低 $D$ が非零で有向三角形を作る条件（定理~\ref{bd20:nonzero-cycle}）に加え，
定理~\ref{bn21:two-negative-obstruction} は負枝を二つ持つ場合を排除する．
既知の全正枝・全負枝の排除と合わせると，負枝 $N$ はちょうど一つ，
正枝 $P_1,P_2$ はちょうど二つでなければならない．

さらに定理~\ref{bs21:two-sectors} により，高 ray を交換して
$\sigma_N=+$ と正規化した node 周期は次の二 sector だけに限られる．
上付き符号は式~\eqref{bn21:node-orientation} の canonical 高 ray 標識 $\sigma$，
$\tau$ は周期に沿う next の低 $D$ の共通符号である．
\begin{center}\small
\begin{tabular}{cclc}
型&$\tau$&node 周期&$(L_{P_1},L_{P_2})$\\\hline
I&$+$&$N^+\to P_1^+\to P_2^+\to N^+$&$(3,3)$\\
II&$+$&$N^+\to P_1^+\to P_2^-\to N^+$&$(4,3)$
\end{tabular}
\end{center}
いずれも負枝を $S_N=T_N=w$，$r_N=w+H$ と書くと
\[
 w\ge135,\qquad H\ge w
\]
が必要である．型 II ではさらに
$S_{P_2}=w-1$，$T_{P_2}=w+1$，辺 $P_2\to N$ の node は $(0,0,5)$ となる．
また系~\ref{ps21:book-denominator-order} により，この正規化の下で
$q_{i0}<q_{i1}$，$i=2,3,4$ が成り立つ．
他の sector の排除は，その存在から解析的に $w\le134$ を導き，
その範囲における全 scale を因数分解で復元した結果である．
二つの残余 sector の実現・排除は未解決であり，book 型は支持表に残る．

\paragraph{低辺補 graph が $P_3$ である $G_{5,15}$．}
低辺を $23,24$，二価 core を $2$ とし，
$Q=q_{34}$，$u=p^3_{24}$，$v=p^4_{23}$ と置く．
定理~\ref{g51521:reduction} は，共通低 puncture の正例外型
$\{u,v\}=\{1,Q-1\}$ を無上界で排除する．
通常型 $v=-u$ では，中心の全 scale は表~\ref{g51521:center-table} の四点に限られ，
三高 cores の低 rows と低 $D$ は表~\ref{g51521:ordinary-row-table} の十四配置に限られる．
負例外型は $Q=3,5$ の二つの template に帰着する
（命題~\ref{g51521:negative-templates}）．従って
\[
 Q\ge7\quad\Longrightarrow\quad
 (q_{20},q_{21})\in\{(7,31),(9,11),(31,7),(19,43)\}
\]
となる canonical 標識を取れる．この結論は $Q$ 自身の上界を与えない．

通常型の十四配置の各々で $Q$ を固定すると，cores $2,3,4$ とその相互の三辺を満たす
整数データを，他の scale や分母の上限なしで有限に全復元できる
（命題~\ref{g51521:fixed-Q-finite}）．
共有高 $D$ の一次式が共鳴する場合も残る pair equation は非退化な二次式である．
全 $Q$ にわたる通常型，負例外型，および四高 cores $0,1$ を加えた接着は未解決であり，
$G_{5,15}$ も支持表から除かれない．

\paragraph{三高係数の非負領域と残る量化．}
三高分母の高係数が全て非負である問題は，
定理~\ref{hd20:nonsingular} により非特異な Schur 領域に限られる．
さらに命題~\ref{hd21:reconstruction} は，指定した高係数
$(x,y,z)=(D_a,D_b,D_c)\ge0$，$x+y+z\ge1$ に対して $x\le s-1$ を要求し，
尺度の整数性を単一合同条件へまとめる．
式~\eqref{hd21:fixed-data} の記号では
\[
 m=F^2+F+1,\qquad m\mid FE-J
\]
であり，$\alpha\ge k\ge b+1$ と元の一つの pair-product 方程式，
および尺度・entry の正値性も課すと，
必要十分な復元判定となる．定理~\ref{hd21:norm-bound} は
\[
 3\le m\le E^2+EJ+J^2,\qquad m\mid E^2+EJ+J^2
\]
を与えるため，五整数 $(s,t,x,y,z)$ ごとに残る $b=-\beta$ の全範囲は有限である．
この五整数自体の一様上限と，同時非負性の一般不存在は未解決である．

$G_{5,15},G_{5,16},G_{5,17}$ は $K_4$ を部分 graph として含む．
四使用頂点 $K_4$ や $K_4+\mathrm{leaf}$ の排除は支持全体を指定した定理であるため，
それだけでこれらの密な graph を除くことはできない．
今回の追加還元は，五型六十六標識という支持の必要候補数を変更しない．
五候補の整数接着・幾何的実現と，一般の五次分母剛性は引き続き未解決である．
